\subsection{Análisis factorial exploratorio}
\label{sec:analisis-factorial}

El análisis factorial complementa al ACP (\cref{sec:pca}) al separar
la varianza compartida entre variables de la específica de cada una.
Para cada variable
estandarizada se plantea
\begin{equation}
  z_j=\sum_{k=1}^{3}\lambda_{jk}f_k+\varepsilon_j,
  \qquad h_j^2=\sum_{k=1}^{3}\lambda_{jk}^2.
  \label{eq:af-modelo}
\end{equation}
Los $f_k$ son factores latentes y las cargas $\lambda_{jk}$ indican
la magnitud y el sentido de su asociación con cada variable. Con
factores ortogonales de varianza unitaria, la comunalidad $h_j^2$ es
la proporción de varianza explicada; la unicidad $1-h_j^2$ recoge
la parte específica y el error.

\paragraph{Adecuación y extracción}
El diagnóstico inicial utiliza 15 predictores y 16,028 observaciones
completas de 2021--2024. Bartlett
rechaza la ausencia de correlaciones ($p<0.001$), aunque la dependencia
temporal y espacial limita esta prueba. El KMO inicial de 0.647
indica adecuación mediocre. Se retira \texttt{NOX\_06\_09} por su
redundancia con NO y NO2 (VIF de 45.73). Sobre las 14 variables
restantes, MAP de Velicer sugiere tres factores, frente a siete del
análisis paralelo; se adopta MAP por parsimonia, conservando esta
discrepancia como limitación.

La extracción por ejes principales iterados produce una comunalidad
inadmisible de 1.093 para \texttt{RH\_min\_diurna}; se elimina esta
variable y se conserva la humedad media. El bloque final de 13
predictores alcanza un KMO de 0.677 y un VIF máximo de 2.54, sin
comunalidades superiores a uno. Se mantienen tres factores y se
aplica varimax para facilitar su interpretación mediante ejes
ortogonales (\cref{fig:af-diagrama}).

\begin{figure}[t]
  \centering
  \includegraphics[width=0.92\linewidth]{af_diagrama_cargas.pdf}
  \caption{Cargas varimax con $|\lambda|\geq0.30$. Flechas sólidas:
  asociaciones positivas; discontinuas: negativas. Los valores grises
  indican unicidades.}
  \label{fig:af-diagrama}
\end{figure}

\paragraph{Interpretación}
Los factores explican el 40.3\,\% de la varianza total de los 13
predictores. \emph{Combustión primaria} (15.7\,\%) agrupa NO2 (0.74), NO
(0.68), PM10 (0.63), CO (0.49) y SO2. \emph{Ventilación} (14.5\,\%) reúne
la velocidad media y mínima del viento, con cargas de 0.88 y 0.77, y la
componente zonal ($-0.52$). \emph{Forzamiento fotoquímico} (10.1\,\%)
combina temperatura máxima (0.72), radiación acumulada (0.63) y humedad
media ($-0.48$): días
cálidos, soleados y secos. La presión tiene comunalidad de apenas
0.13, por lo que el nombre ventilación resulta más preciso que
atribuir al factor un mecanismo anticiclónico.

Tres decisiones que se probaron y se descartaron en el camino ---el análisis paralelo de Horn, la extracción por máxima verosimilitud y la rotación oblicua--- se documentan en \cref{sec:tecnicas-descartadas}.

\paragraph{Ajuste y alcance}
El residuo cuadrático medio de las correlaciones reconstruidas
(RMSR) es 0.0555; 26 de 78 pares conservan residuos absolutos mayores
que 0.05. El ajuste es razonable, pero parcial. Los puntajes de
Bartlett, aplicados a 2025 con pesos de entrenamiento, tienen VIF
cercanos a 1.01 y correlación absoluta máxima de 0.097. Por eso son
los predictores de los modelos de \cref{sec:clasificacion_tecs}, en lugar
de las puntuaciones del ACP, cuyo contraste se reporta ahí.

Los tres factores reaparecen en 2025 con congruencia de Tucker de
0.973 para combustión, 0.939 para ventilación y 0.924 para fotoquímica:
la estructura es estable fuera del periodo de ajuste. La determinación de
los puntajes (correlación entre puntaje y factor) es 0.880, 0.918 y 0.833;
el forzamiento fotoquímico comparte 69\,\% de su varianza con el factor
que representa, y el resto es error de estimación. Persisten el KMO
mediocre y la transformación Box--Cox estimada sobre todo el periodo. Los factores describen asociaciones; sin mediciones de
compuestos orgánicos volátiles no permiten establecer mecanismos
de formación de ozono.
